\section{A complete analytical family with A delay and active renewal}
\label{sec:witness}
This consolidates the OR-Q1 construction with W-F at its fixed anchor
$v_0=v_+$. There the old trial and selected work expressions coincide
exactly. The result is coexistence under the retained equations, not a
calibrated scenario, numerical rollout, chosen task setting or optimal-
advice certificate. Long geometry and very slow recovery may be needed.

\subsection{H1 input and finite-prefix auxiliary limit}
Translate the eligible front to time zero. A donor-limited H1 pulse at
$v_0$ has arrival rates
\begin{equation}
 e(t)=e_f e^{-kt},\quad a(t)=s-e(t),\quad
 k=v_0/h>0,\quad s=v_0r<R,\quad\varepsilon=R-s>0.
 \label{eq:winput}
\end{equation}
Only in this construction, $e,a$ denote arrival rates, not densities;
$e_f/v_0$ is the corresponding front density. Choose $c_0>1$, $e_f>0$,
$a_f=s-e_f>0$ with
\begin{equation}
 \varepsilon<\delta e_f/2,\quad
 G_0=(\varepsilon+e_f)/c_0<R,\quad
 d_0=R-G_0>\varepsilon+e_f,\quad\delta=c_0-1.
 \label{eq:wconditions}
\end{equation}
The initial free work is $a_f+c_0e_f=R-\varepsilon+\delta e_f>R$.
E therefore queues first under the selected priority.

Temporarily set $a_r=\beta=0$ solely as an auxiliary finite-prefix limit.
Before window fill, while A is free and E queued,
\begin{equation}
 g_E=\frac{\budget-a(t)}{c_0},\quad
 \dot\budget=-\frac{g_E(R-g_E)}m,\quad
 \dot q_E=e-g_E. \label{eq:waux}
\end{equation}
This follows from the zero-aging profile and
$D(0;U(g_E))=R-g_E$. Since $g_E\leq G_0$ and
$\budget\geq a\geq a_f$,
\begin{equation}
 R-\budget\geq d_0Y_E/m,\qquad
 Y_E\leq m(\varepsilon+e_f)/d_0<m. \label{eq:wprefill}
\end{equation}
A first fill before A onset would contradict this bound, so the unfilled
prefix is self-consistent.

While $\budget\geq s$ and $e\geq e_f/2$,
\begin{equation}
 -\dot\budget\geq\mu:=\frac{e_fd_0}{2mc_0}>0,\qquad
 \dot q_E\geq(\delta e_f/2-\varepsilon)/c_0>0.
 \label{eq:wdrop}
\end{equation}
Take positive $k$ sufficiently small that
$k\varepsilon/\mu<\log2$. The budget reaches $s$ by
$\varepsilon/\mu$ while the condition on $e$ holds, then passes below
$s$. Afterwards, while A is still free, $g_E<e/c_0<e$, preventing E
clearance. Because $a(t)\uparrow s$ and the budget is nonincreasing, a
finite first A onset $t_A$ occurs. Its pre-hit margin $x=\budget-a$ has
\begin{equation}
 x(t_A)=0,\quad\dot x(t_A-)=-ke(t_A)<0,\qquad
 \dot q_A(t_A+)=\budget(t_A)/2>0. \label{eq:wAbirth}
\end{equation}
Thus increased A arrivals can cross the protected-service margin even
though the constant-A limit in \eqref{eq:constantA} cannot.

\subsection{A backlog large enough to exercise renewal}
At $k=0$ in the same auxiliary input, A is constant $a_f$. The finite-time
barrier keeps $\budget>a_f$, $Y_E<m$ and E queued, with
\begin{equation}
 \dot q_E\geq\sigma:=(\delta e_f-\varepsilon)/c_0>0.
 \label{eq:wbacklog}
\end{equation}
Choose finite $T_b$ with $\sigma T_b>3m$. The free-A margin remains
strict on that compact interval. Continuous dependence of
\eqref{eq:waux} on its input permits a sufficiently small positive $k$
for A to remain free through $T_b$ with $q_E(T_b)>2m$. Retain
\eqref{eq:wdrop}. The finite A onset is later and E cannot clear in between.
This removes constant input as a primary assumption; $k=0$ is only a
finite-prefix comparison.

End the actual finite pulse at $T>t_A$ sufficiently close to $t_A$ that
$q_A(T)>0$, $q_E(T)>2m$, $Y_E(T)<m$, $\budget(T)>0$, and
\begin{equation}
 \frac{q_A(T)}{c_0}<m-Y_E(T),\quad
 q_A(T)<c_0q_E(T),\quad
 \frac{Rq_A(T)}{mc_0}<\budget(T)/2. \label{eq:wsmallA}
\end{equation}
The first A queue tends to zero as $T\downarrow t_A$, while the other
margins stay positive, so these inequalities are attainable.

In the auxiliary post-pulse continuation, both queues drain with
$g_A=\budget/2$, $g_E=\budget/(2c_0)$ before fill. A clears first,
using $q_A(T)/c_0$ E mass. Since
$d\budget/dY_E=-(R-g_E)/m\geq-R/m$, \eqref{eq:wsmallA} keeps
$\budget\geq\budget(T)/2$ until this strict clearance, before fill.
Afterwards A has no arrivals and $g_E=\budget/c_0$.

The first $Y_E=m$ is reached in finite time with E still queued. Indeed,
before it,
\begin{equation}
 \budget=R(1-Y_E/m)+\frac1m\int_0^{Y_E}g_E(T(\xi))\dd\xi.
 \label{eq:wfillidentity}
\end{equation}
If $Y_E$ only approached $m$ asymptotically, the positive earlier integral
would give a positive limiting budget and hence positive limiting service,
a contradiction. At fill, $g_E>0$ and the oldest profile has positive
loss. Thus the retirement term in \eqref{eq:renewalupdate} is nonzero
immediately afterwards. The construction exercises genuine renewal, not
only its initial additive prefix.

\subsection{Removing the auxiliary zero parameters}
On a fixed finite interval the positive-recovery kernel converges uniformly
to $\pos{R-c(u)}$ as $a_r\downarrow0$, with error at most $a_rk_j$
times interval length. The root is uniformly regular for sufficiently
small $\beta$ satisfying \eqref{eq:workdomain}. The Volterra estimates
give finite-prefix continuous dependence on these parameters in each
valid regime.

The constructed sequence has strict margins or transverse hits: E front
shortfall, A onset \eqref{eq:wAbirth}, pulse end with positive queues,
A clearance with positive service, and fill with $g_E>0$. Local continuous
dependence and the nonzero guard derivatives preserve the finite sequence
and its inequalities for sufficiently small \emph{positive} $a_r$ and
\emph{positive} $\beta$. A has positive waiting area and oldest-profile
retirement remains nontrivial. The auxiliary zero-recovery limit is not
part of the admitted model; its infinite-time behavior is not used.

For orientation, an all-history additive O-S account started identically
agrees with O-R until first fill $t_R$. At that common state only O-R
retires old weight, so
\begin{equation}
 (\dot\budget_{O\text{-}R}-\dot\budget_{O\text{-}S})(t_R+)
 =\frac{g_E(t_R)}mD(t_R;u_E(0+))>0. \label{eq:wrival}
\end{equation}
This is a local difference, not global policy ordering or evidence that
one closure is physically correct. O-S remains unselected.

After the pulse, no new queues are born. At most two clearances remain;
the positive-budget/clearance bounds make them finite. The small A tail
preserves A-first clearance, then E alone clears. This gives a complete
unique regular family with positive A delay and active renewal.

\subsection{Realizing it with the unchanged H1 geometry}
Choose $r=s/v_0$, $e_0+a_0=r$, $a_0=b_0>0$. A receiving-limited
initial branch followed by a donor-limited front is possible when
\begin{equation}
 \kc(v_0)-r<(\gamma-1)e_0,\qquad
 e_f/v_0<\frac{\kc(v_0)-r}{\gamma-1}. \label{eq:wh1}
\end{equation}
Choose front progress $z_f$ after the switch. Increasing $h$ at fixed
$z_f/h$ retains that density while making $k=v_0/h$ small. Put
\begin{equation}
 \ell=v_0T,\quad L_p=\ell+z_f,\quad
 M=\ell(r+b_0),\quad b>M/k_j,\quad L=L_p+b. \label{eq:wgeometry}
\end{equation}
The terminal length changes only the common translated eligibility delay,
not the node argument. All original advice-domain and far-lane conditions
can be retained; no advice subinterval is silently substituted.

An exact normalized static feasibility point from the earlier proof is
\begin{equation}
 R=m=1,\quad c_0=11/10,\quad e_f=1/8,\quad
 \varepsilon=1/512,\quad s=511/512, \label{eq:wstatic}
\end{equation}
giving
\begin{equation}
 \begin{gathered}
 d_0=2491/2816,\quad\mu=12455/247808,\quad
 \varepsilon/\mu=484/12455,\\
 \sigma=27/2816,\qquad (\varepsilon+e_f)/d_0=715/4982<1.
 \end{gathered} \label{eq:wstaticderived}
\end{equation}
One compatible H1 choice uses
\begin{equation}
 \begin{gathered}
 v_-=1,\quad v_0=v_+=2,\quad w=2,\quad k_j=3,\\
 a_0=b_0=1/8,\quad e_0=383/1024,\quad
 r=511/1024,\quad\gamma=5,\quad R_B=1.
 \end{gathered}
 \label{eq:wstaticH1}
\end{equation}
Its fast front density $1/16$ is reached at
\begin{equation}
 z_f/h=5\log(5632/5125)+\log(1025/256). \label{eq:wfront}
\end{equation}
These are static algebraic existence checks, not selected traffic
parameters. $k$, terminal geometry and positive recovery/mismatch values
are chosen by strict finite-prefix arguments, not by a numerical rollout.
The witness persists locally near the anchor; the separate forward theorem
qualifies all original advices, without claiming A delay at all of them.
No preferred advice, interior optimum or useful gradient is established.

\section{What physical observations can and cannot identify}
\label{sec:physical}
This collects MP6-Q1's conditional algebra; it neither executes nor opens
calibration. The required informative observations may not be available in
the primary cohort. A future execution proposal must specify their scale,
uncertainty, tolerances and development/held-out separation.

\subsection{Recovery tails and immediate work}
If no more E is entering and A remains queued, $g_E=0$, $c_A=1$, and
work conservation gives
\begin{equation}
 g_A=\budget=R-\loss,\qquad\loss=R-g_A. \label{eq:tailidentify}
\end{equation}
With independently known $R$, the ideal saturated tail exposes the account
without current E maneuver work. A free A tail does not: low outflow may
simply be low demand. One tail need not identify every profile parameter
or test overlapping-history retirement.

At a binding mixed observation with $g_E>0$, set
\begin{equation}
 X=g_A\pos{u_A-u_E},\qquad
 Z=\frac{R-\loss-g_A}{g_E}-1=\delta+\chi X. \label{eq:workidentify}
\end{equation}
Once recovery and marks are separately constrained, two distinct $X$
values give
\begin{equation}
 \chi=\frac{Z_2-Z_1}{X_2-X_1},\qquad
 \delta=\frac{X_2Z_1-X_1Z_2}{X_2-X_1}. \label{eq:rankidentify}
\end{equation}
The design determinant is $X_2-X_1$. Equal covariates do not separate
the coefficients; all $X=0$ leaves $\chi$ unresolved. Nonbinding service
gives only a feasibility inequality. Admissibility still requires the
nonnegative and root-domain bounds, not merely a fitted equality.

Without independent recovery information, a single observation has the
accounting ambiguity
\begin{equation}
 \delta'=\delta+\epsilon,\quad
 \loss'=\loss-\epsilon g_E,\quad
 \loss'+g_A+(1+\delta'+\chi X)g_E
 =\loss+g_A+(1+\delta+\chi X)g_E. \label{eq:confounding}
\end{equation}
This is not an alternative O-R solution or a global identifiability theorem:
the primed account need not satisfy the history law. It shows why a total
discharge/TST fit alone cannot establish that present work and retained
restriction describe independent real effects. Adding $W$ again to $c_E$
would double-charge the modeled profile loss. Distinct empirical attribution
still requires tails, mixtures and joint predictions with one parameter set.

\subsection{Preparation and stationary longitudinal correspondence}
From \eqref{eq:h1donor}--\eqref{eq:h1receiving}, informative pure branches
give
\begin{equation}
 -\frac{d\log e}{dz}=1/h\quad\hbox{(donor)},\qquad
 -\frac{d\log(e+\Delta)}{dz}=1/(\gamma h)\quad\hbox{(receiving)}.
 \label{eq:prepidentify}
\end{equation}
The receiving branch alone primarily identifies the product. The
transformation $h'=sh$, $\gamma'=\gamma/s$ preserves it provided the
branch and domain remain unchanged. A donor branch can constrain $h$
separately; a final E count generally cannot identify the full law.

For a documented ideal constant-speed headway relation with standing
spacing $d$ and desired time gap $\tau_d$,
\begin{equation}
 k_{\rm ideal}(s)=\frac1{d+\tau_ds},\qquad
 f_{\rm ideal}(s)=\frac{s}{d+\tau_ds}. \label{eq:headway}
\end{equation}
Setting $k_j=1/d$ and $w=d/\tau_d$ makes
$\kc(s)=k_{\rm ideal}(s)$ and $c(s)=f_{\rm ideal}(s)$. This is a
stationary curve identity at actual speed $s$. It establishes neither
native transient wave speed nor instantaneous advice compliance, and it
does not prove the stream-specific queued mark $U(g_j)$ at a merge.

\subsection{Joint physical admission remains unresolved}
One common physical parameter set must explain upstream transfer, actual
crossing-speed marks, mixed discharge and recovery while satisfying the
full-advice primary domain. Actual acceleration constrains the source
meaning of $a_r$; discharge tests its boundary effect. $V$ must not be
fitted as a vehicle-speed trajectory. The physical renewal normalization
is fixed; a different window would require reopening its contract.

Development trajectory observations are evidence for common parameters,
not extra future inputs to a controller or prescribed native outcomes.
Fit neither preferred actions nor gradient signs, and do not select only
interior-optimum cases. The mathematical witness and illustration values
from source papers are not a jointly supported applied range. No
theory-only primary is adopted; that would need an explicit role decision.

\section{Provenance, artifact map and remaining scientific boundary}
\label{sec:provenance}
The following source records were inspected for consolidation. They retain
their original equations and delivery-time statuses. This report assembles
only the selected current member and labels historical results explicitly.

\begin{longtable}{>{\raggedright\arraybackslash}p{35mm}>{\raggedright\arraybackslash}p{110mm}}
\toprule
Mathematical responsibility & Archived record and use here\\
\midrule
\endhead
Original S13 &
\href{../s13/s13_m2_model_contract.pdf}{Model contract} and
\href{../s13/s13_m2_certificate_analysis.pdf}{certificate analysis};
Section~\ref{sec:s13}, unchanged historical reference.\\[6pt]
Preparation, footprint, finite terminal &
\href{vsl_h1_t1_joined_contract.pdf}{H1/T1 joined contract};
Sections~\ref{sec:domain}--\ref{sec:terminal}. Its old strict-priority
one-queue server is superseded, not included as current service.\\[6pt]
Root, priority and clearance convention &
\href{vsl_se_admission_and_transition_audit.pdf}{S-E audit} and
\href{vsl_sep_recovery_and_clearance_audit.pdf}{S-E-P/recovery audit};
Sections~\ref{sec:root}--\ref{sec:allocation}, with subsequently selected
W-F/R-E/O-R and global qualification.\\[6pt]
Recovery update, causal bound and witness &
\href{vsl_or_recovery_qualification.pdf}{O-R qualification};
Sections~\ref{sec:recovery}, \ref{sec:local}, \ref{sec:witness}.
Its advice-normalized trial work is replaced by W-F; the witness is
preserved at the exact anchor.\\[6pt]
Fixed coefficient and advice causality &
\href{vsl_j4_speed_response_audit.pdf}{J4 speed-response audit} and
\href{../../reporting/research_readiness/vsl_wf_selection_and_scope.md}{W-F selection};
Section~\ref{sec:root}. W-A is unselected.\\[6pt]
Whole original domain, events, clearance &
\href{vsl_wf_forward_qualification.pdf}{WF-Q1 forward qualification};
Sections~\ref{sec:global}--\ref{sec:information}.\\[6pt]
Physical interpretation and identification &
\href{vsl_mp6_physical_identifiability.pdf}{MP6-Q1 note} and
\href{vsl_mp6_joint_physical_admission_contract.md}{joint admission contract};
Section~\ref{sec:physical}. Applied admission remains open.\\
\bottomrule
\end{longtable}

The already reviewed primary literature motivates separate ingredients:
lane conservation and mandatory transfer; receiving-work and allocation;
boundary-count propagation; and bounded-acceleration point insertion.
The exact homogeneous preparation, stream speed/work closure, causal
priority and recent-mass renewal remain our adaptations. Their proof
arguments are not attributed to a generic CTM/LTM theorem or a published
complete lane-drop model.

\begin{center}
\href{../../extraction/research_readiness/vsl_feedback8_model_stack_source_map.md}{Detailed archived source and decision map}
\end{center}

\textbf{What is complete.} All conditional forward equations, initialization,
units, selected mode rules, history update, conservation and objective are
specified. Existing mathematical arguments give unique continuation and
finite clearance over the retained original advice interval, as well as
a nonempty analytical coupled-delay family. The symbolic verification
checks supporting identities and rational static inequalities only.

\textbf{What is not established.} No calibrated joint regime, useful action
split, full-lane/discrete-order necessity, AD correctness, empirical recovery
law or external policy advantage is claimed. Local activation-flux gaps
do not prove a full-objective jump. Native finite-cohort requeueing is not
automatically invalid. Strong aggregate, scalar and search rivals remain.

Feedback 8 is now answered by consolidation, not by executing the next
protocol block. The active checkpoint remains MP6-Q1. Applied MP6 and
Step 3 stay open; Step 4 and simulator/experiment work remain unopened.
